\subsection{问题一分析}
问题一同时考察"鱼类资源量增加"、"鱼类增长的驱动机制"和"基础资源的承载压力变化"。若把四大家鱼与水草、浮游植物、浮游动物、底栖饵料各自合并，模型仍能匹配2025年1.8倍的总量，却无法分辨哪一层资源先承压，也难以解释"水下荒漠"等底层退化现象。因此，方程中的四类资源必须与四大家鱼的食性逐项对应。

题目给出的两类观测数据分别对应两种尺度：2025年四大家鱼1.8倍恢复量对应全流域平均口径，可用于约束系统总量参数；湖北江段单网次渔获量提升120%具有明显局地性和网具特异性，更适合作为外部独立验证，不宜与全流域总量数据共同参与参数标定。据此，我们先在统一总量尺度下识别禁渔的直接生态效应，再将局地捕获差异纳入结果分析。

